% Source: https://artofproblemsolving.com/community/c3h1967383_intro_to_algebra
% Stable identity: aops-topic-1967383-q01
Multiply by $mn$: $4n+2m=mn$, so $mn-2m-4n=0$ and
\[
(m-4)(n-2)=8 .
\]
Since $\dfrac{4}{m}<1$ and $\dfrac{2}{n}<1$, we have $m>4$ and $n>2$, so both factors are positive. The factor pairs of $8$ are $(1,8),(2,4),(4,2),(8,1)$, giving
\[
\boxed{(m,n)=(5,10),\ (6,6),\ (8,4),\ (12,3)} .
\]
So yes, $(12,3)$ and $(6,6)$ are solutions: $\dfrac{4}{12}+\dfrac{2}{3}=1$ and $\dfrac{4}{6}+\dfrac{2}{6}=1$. The list $(8,1),(1,8),(4,2),(2,4)$ quoted in the post is the list of factor pairs $(m-4,n-2)$, not of the pairs $(m,n)$ themselves.
